% =====================================================================
%  Secao para a monografia do TCC -- fragmento para \input{} no documento
%  principal (NAO tem \documentclass nem \begin{document}).
%
%  Depende apenas de comandos ja presentes no preambulo do documento
%  principal:
%    - listings + \lstdefinelanguage{mermaid} + \lstset (frame=single)
%    - amsmath                (equacao do PISA)
%    - todonotes + \fabio{}/\guilherme{}  (comentarios de revisao)
%    - hyperref               (\url)
%    - babel[brazil], booktabs
%
%  Observacao: o texto DENTRO dos blocos lstlisting[language=mermaid] e
%  mantido em ASCII (sem acentos) porque o pacote listings, do jeito que
%  esta configurado no preambulo, nao mapeia caracteres UTF-8 acentuados.
%  Para usar acentos nos diagramas, adicione um "literate" ao \lstset.
%  O texto corrido (fora dos diagramas) usa acentuacao normal.
%
%  Fonte tecnica: docs/ARCHITECTURE.md, PROTOCOL_SPEC.md, FIRMWARE.md e
%  docs/BLE_PAYLOAD_VALIDATION.md do repositorio do simulador.
% =====================================================================

\section{Arquitetura e fluxos de interação do simulador}
\label{sec:arquitetura-fluxos}

O simulador é composto por duas metades que se comunicam exclusivamente por
Bluetooth Low Energy. No computador, a aplicação de mesa (Python/PyQt6) é
responsável pela interface gráfica, pelas sessões BLE, pelo contrato de
formato de mensagens e por uma segunda cópia, independente, da matemática dos
modelos fisiológicos, executada localmente para produzir uma curva de
referência sem ruído. Na placa de desenvolvimento (C sobre o Zephyr RTOS, por
meio do nRF Connect SDK), uma \textit{thread} avança o modelo fisiológico e o
modelo de ruído do sensor uma vez por segundo de relógio real, enquanto uma
segunda \textit{thread} é responsável pelo envio BLE e por drenar uma fila de
escritas de configuração, persistindo cada uma em memória \textit{flash}
SPI-NOR externa. A placa é autônoma: desconectada da aplicação, continua
anunciando e transmitindo.

Todo o tráfego que cruza o enlace BLE pertence a uma de três classes:
(i)~leituras de glicose no perfil padrão da Bluetooth SIG (\textit{Continuous
Glucose Monitoring Service}, UUID \texttt{0x181F}), da placa para a aplicação;
(ii)~configuração do simulador sobre um serviço GATT proprietário (base
\texttt{5b2c0001-\ldots}), nos dois sentidos; e (iii)~uma notificação de um
\textit{byte}, chamada \textit{reset-sync}, da placa para a aplicação, cujo
\emph{instante de chegada} sinaliza que uma escrita de configuração fez
efeito. As duas simulações (placa e aplicação) rodam em relógios separados,
re-ancorados em $t = 0$ apenas quando uma escrita de configuração reinicia
ambas; o ruído do sensor, presente somente na placa, é a única diferença que
deve separar visivelmente as linhas ``recebida'' e ``esperada'' no gráfico.

\subsection{Fluxos de interação}
\label{sec:fluxos}

Cada fluxo a seguir é apresentado como um diagrama de sequência (na notação
\textit{mermaid}) acompanhado da explicação do seu propósito. Os atores são:
\textbf{Usuário}; \textbf{App}, as janelas PyQt; \textbf{BleSession}, uma
conexão BLE persistente com laço \textit{asyncio} próprio; e as duas
\textit{threads} do firmware, \texttt{comm\_thread} e \texttt{model\_thread}.

\subsubsection{Fluxo de desenvolvimento: início do streaming}
\label{sec:fluxo-dev}

É o caminho de uso mais simples: placa e aplicação já estão configuradas (por
valores padrão ou por um envio anterior) e o objetivo é apenas observar as
duas curvas. Ao clicar em \emph{Start}, a aplicação ancora o seu motor de
simulação local em $t = 0$ e escreve a característica de estado de execução;
a partir daí cada lado avança de forma independente, no seu próprio
temporizador. A placa envia, a cada \texttt{measurement\_interval} segundos
(cinco, nesta versão), uma notificação de \textit{CGM Measurement} e uma de
\textit{Food/Exercise Status}; a aplicação anexa a primeira à linha
``recebida'' e a segunda ao gráfico de refeição e exercício, enquanto o motor
local alimenta a linha ``esperada''.

\begin{lstlisting}[language=mermaid,caption={Fluxo de desenvolvimento: do clique em Start ao streaming contínuo.},label={lst:fluxo-dev}]
sequenceDiagram
    participant U as Usuario
    participant App as App (graphic/)
    participant BLE as services/ble_session.py
    participant FW as Firmware (comm_thread)
    participant Model as Firmware (model_thread)

    U->>App: Start
    App->>App: ancora SimulationEngine local em t=0
    App->>BLE: queue_write("run_state", RUNNING)
    BLE->>FW: escrita GATT
    FW->>Model: apply_config()
    loop a cada 1 s (tick do MCU)
        Model->>Model: passo do modelo + ruido do sensor
    end
    loop a cada 5 s (measurement_interval)
        FW->>BLE: notify CGM Measurement (glicose)
        FW->>BLE: notify Food/Exercise Status (taxa de carbo, % exercicio)
        BLE-->>App: sinal new_message
        App->>App: anexa a linha "recebida"
    end
    App->>App: engine local avanca a cada 1 s -> linha "esperada"
\end{lstlisting}

\subsubsection{Fluxo de envio de configuração}
\label{sec:fluxo-envio}

É acionado pelo botão \emph{Enviar para a placa} de cada janela de
configuração (Pessoa, Sensor, Refeição, Exercício). A escrita sempre provoca
um reinício completo na placa (a rotina \texttt{apply\_config\_locked()}
re-inicializa todos os \textit{slots} e zera \texttt{sim\_clock\_min}), e é
justamente assim que a aplicação confirma que a escrita chegou: ela aguarda a
próxima notificação \textit{reset-sync} e exibe ``Aplicado na placa''. Em uma
placa multi-sensor, a escrita atinge o \textit{slot} apontado pelo cursor
\textit{Sensor select} (característica \texttt{5b2c0015}); por isso, quando um
\textit{slot} específico é o alvo, a aplicação enfileira antes uma escrita
nesse cursor.

\begin{lstlisting}[language=mermaid,caption={Fluxo de envio de configuração, com o reinício e a confirmação por reset-sync.},label={lst:fluxo-envio}]
sequenceDiagram
    participant U as Usuario
    participant App as App (*_config_window.py)
    participant BLE as services/ble_session.py
    participant FW as Firmware (comm_thread)
    participant Model as Firmware (model_thread)

    U->>App: edita e clica "Enviar para a placa"
    opt alvo de um slot especifico (N > 1)
        App->>BLE: queue_write("sensor_select", slot)
        BLE->>FW: escrita GATT -- comm_thread ajusta working_sel
    end
    App->>BLE: queue_write(char_key, bytes)
    BLE->>FW: escrita GATT (caracteristica de config)
    FW->>FW: atualiza sim_config na RAM e grava na flash externa
    FW->>Model: apply_config() -- reinicia todo slot, sim_clock_min = 0
    Model-->>FW: notify reset_sync
    FW-->>BLE: notificacao
    BLE-->>App: sinal reset_sync
    App-->>U: "Aplicado na placa"
\end{lstlisting}

Uma escrita na característica de \emph{Speed} (\texttt{5b2c0012}) hoje passa
por esse mesmo caminho de reinício, zerando o relógio e limpando os vetores
de eventos pontuais. Mudar a velocidade de reprodução não deveria perturbar
uma simulação em andamento; tratar a velocidade como um escalar aplicado ``a
quente'' é uma correção pendente.
\guilherme{Decidir se essa limitação de Speed entra como ``trabalho futuro'' ou como nota de implementação.}

\subsubsection{Fluxo de leitura de configuração}
\label{sec:fluxo-leitura}

É acionado pelo botão \emph{Ler da placa}: uma operação de leitura e resposta
GATT simples, sem notificação e sem reinício. O pedido de leitura trafega na
\emph{mesma fila} das escritas dentro da \texttt{BleSession}, de modo que a
sequência ``ajustar o cursor \textit{Sensor select} e então ler aquele
\textit{slot}'' permanece ordenada mesmo enquanto a \texttt{comm\_thread}
está ocupada enviando medições.

\begin{lstlisting}[language=mermaid,caption={Fluxo de leitura de configuração (leitura/resposta GATT, sem reinício).},label={lst:fluxo-leitura}]
sequenceDiagram
    participant U as Usuario
    participant App as App (*_config_window.py)
    participant BLE as services/ble_session.py
    participant FW as Firmware (comm_thread)

    U->>App: clica "Ler da placa"
    opt alvo de um slot especifico (N > 1)
        App->>BLE: queue_write("sensor_select", slot)
        BLE->>FW: escrita GATT
    end
    App->>BLE: request_read(char_key)
    BLE->>FW: leitura GATT (caracteristica de config)
    FW-->>BLE: bytes crus (sim_config na RAM)
    BLE-->>App: config_read(address, char_key, raw_bytes)
    App->>App: protocol.decode_*(raw_bytes)
    App->>App: carrega no formulario
\end{lstlisting}

\subsubsection{Fluxo de eventos pontuais (refeição, exercício e PISA)}
\label{sec:fluxo-pontual}

Os comandos ``Inserir Refeição/Exercício/PISA agora'' deliberadamente
\emph{não} seguem o fluxo de reinício. O evento único é acrescentado a um
pequeno conjunto fixo de posições de evento ativo dentro da
\texttt{model\_thread} e decai sozinho ao longo da sua duração, podendo ser
injetado no meio de uma simulação sem reiniciar o relógio nem o estado do
modelo --- por isso não há troca de \textit{reset-sync}. Em uma placa
multi-sensor, a aplicação escreve antes o cursor \textit{Sensor select} e
envia o evento a \emph{uma} sessão (e não em \textit{broadcast} para as
quatro, o que aplicaria o evento $N$ vezes). O motor local da aplicação
espelha o mesmo evento, para que a linha ``esperada'' se mova de forma
idêntica.

\begin{lstlisting}[language=mermaid,caption={Fluxo de eventos pontuais: injeção sem reinício em uma simulação em andamento.},label={lst:fluxo-pontual}]
sequenceDiagram
    participant U as Usuario
    participant App as App (graphic/)
    participant BLE as services/ble_session.py
    participant FW as Firmware (comm_thread)
    participant Model as Firmware (model_thread)

    U->>App: Inserir Refeicao/Exercicio/PISA agora
    App->>App: SimulationEngine.add_instant_*() (linha "esperada" local)
    opt multi-sensor: alvo de um slot
        App->>BLE: queue_write("sensor_select", slot)
    end
    App->>BLE: queue_write("*_instant", bytes)
    BLE->>FW: escrita GATT (one-shot)
    FW->>Model: model_thread_add_instant_*()
    Note over Model: decai ao longo de duration_min, somado/maximizado com a agenda a cada tick -- sem reinicio, sem reset_sync
\end{lstlisting}

Um reinício \emph{posterior}, vindo de qualquer outra escrita (inclusive
\emph{Stop}), limpa esses eventos junto com todo o restante do estado.

\subsubsection{Fluxo de envio de \textit{layout} da placa (multi-sensor)}
\label{sec:fluxo-layout}

A janela \textbf{Board Layout} associa um perfil salvo de Pessoa e de Sensor a
cada um dos $N$ \textit{slots} e envia o conjunto em uma única ação. O método
\texttt{BleSession.send\_board\_layout()} executa a operação como uma única
co-rotina sobre \emph{uma} conexão (qualquer identidade alcança o serviço de
configuração compartilhado), com ritmo controlado para a fila do firmware
acompanhar, subindo uma trilha CSV por \textit{slot} quando a pessoa
atribuída é baseada em CSV e terminando com uma escrita de estado de execução
igual a RUNNING.

\begin{lstlisting}[language=mermaid,caption={Fluxo de envio de layout multi-sensor: um slot por vez, sobre uma única conexão.},label={lst:fluxo-layout}]
sequenceDiagram
    participant U as Usuario
    participant BLW as graphic/board_layout_window.py
    participant BLE as services/ble_session.py
    participant FW as Firmware (comm_thread)
    participant Model as Firmware (model_thread)

    U->>BLW: atribui slots 0..N-1 -> pessoa / CSV, clica "Enviar layout para a placa"
    BLW->>BLW: _build_slots() -- codifica person/sensor/data_source/food/exercise por slot
    BLW->>BLE: send_board_layout(slots)
    loop para cada slot atribuido i
        BLE->>FW: escreve sensor_select = i
        BLE->>FW: escreve person, sensor, data_source, food, exercise (com ritmo, re-tentativa em fila cheia)
        opt pessoa baseada em CSV
            BLE->>FW: CSV BEGIN / DATA x k / COMMIT (na trilha de flash do slot i)
        end
        BLE-->>BLW: board_layout_progress(i+1, N)
    end
    BLE->>FW: escreve run_state = RUNNING
    FW->>Model: apply_config() por escrita -- cada slot reinicia, sim_clock_min = 0
    BLE-->>BLW: board_layout_finished(ok)
\end{lstlisting}

O mapa \textit{slot}~$\rightarrow$~paciente é persistido no lado da aplicação
em \texttt{data/board\_layout.json}. Uma vez atribuído um \textit{slot}, a
lista de dispositivos Bluetooth, a linha da árvore e o título do gráfico
passam a mostrar o nome do paciente no lugar de ``Nordic Glucose Sensor~$N$''.

\subsubsection{Fluxo do modo CGMS Only}
\label{sec:fluxo-cgms-only}

É o oposto do modo ``Model Only''. Com \emph{CGMS Only} marcado, a placa se
comporta como uma fonte de dados puramente CGM-padrão: transmite apenas
notificações de \textit{CGM Measurement}, suprime as de \textit{Food/Exercise
Status} e rejeita toda escrita de configuração, exceto o próprio estado de
execução e a característica CGMS Only. Ativar o modo nunca reinicia (o que
estiver rodando continua); desativá-lo executa o mesmo reinício completo de
uma parada e então aguarda, exigindo um \emph{Start} explícito em seguida ---
diferente do comportamento autônomo padrão da placa, que retoma sozinha.

\begin{lstlisting}[language=mermaid,caption={Fluxo do modo CGMS Only: a placa como fonte CGM-padrão isolada.},label={lst:fluxo-cgms-only}]
sequenceDiagram
    participant U as Usuario
    participant App as App (graphic/main_window.py)
    participant BLE as services/ble_session.py
    participant FW as Firmware (comm_thread)
    participant Model as Firmware (model_thread)

    U->>App: marca "CGMS Only"
    App->>App: trava controles de config e Start/Pause/Stop; descarta SimulationEngine local
    App->>BLE: queue_write("cgms_only", 1)
    BLE->>FW: escrita GATT
    FW->>Model: set_run_state(RUNNING) se ocioso -- sem reinicio
    FW->>FW: ativa rejeicao de escrita nas caracteristicas de config
    loop a cada 5 s
        Model->>Model: passo do modelo + ruido do sensor
        FW->>BLE: apenas notify CGM Measurement (Food/Exercise Status suprimido)
        BLE-->>App: sinal new_message
        App->>App: anexa a linha "recebida" (sem linha "esperada")
    end
    U->>App: desmarca "CGMS Only"
    App->>BLE: queue_write("cgms_only", 0)
    BLE->>FW: escrita GATT
    FW->>Model: set_run_state(STOPPED) -- reinicio completo, depois aguarda
    FW->>FW: desativa rejeicao de escrita
    App->>App: destrava controles
    Note over U,App: placa aguarda Start explicito -- sem retomada automatica
\end{lstlisting}

\subsection{Perfil de comunicação no estilo Dexcom}
\label{sec:perfil-dexcom}

O perfil padrão e principal da placa é o \textit{Continuous Glucose
Monitoring Service} da Bluetooth SIG (UUID \texttt{0x181F}), um padrão
público que qualquer cliente CGM genérico consegue ler. Além dele, a placa
oferece um perfil opcional que \emph{imita} um transmissor Dexcom de
consumidor, selecionável em tempo de execução pela característica de perfil
de comunicação (\texttt{5b2c0014}) e disponível apenas em \textit{builds} de
sensor único. Esse perfil não resulta de engenharia reversa realizada neste
trabalho --- atividade explicitamente fora do escopo ---, mas do uso de
informação já publicada por projetos de código aberto da comunidade de
diabetes (DiaBLE, xDrip+, LoopKit). Dessas fontes sabe-se que sensores Dexcom
reais não implementam o serviço da SIG: eles expõem um serviço GATT
proprietário (UUID curto \texttt{FEBC}), com uma etapa de
autenticação por desafio e resposta (J-PAKE nos modelos G7, AES nos G6) antes
de qualquer dado de glicose, e mensagens binárias marcadas por \textit{opcode}
transmitidas em cadência fixa de 300\,s. Não existe especificação oficial
pública da Dexcom, e por isso nenhuma é citada.

O que o firmware efetivamente implementa é um esboço mínimo, e não um clone
compatível byte a byte: o serviço \texttt{FEBC} com uma característica de
controle cujas escritas são confirmadas e ignoradas (sem qualquer
\textit{handshake} de autenticação) e uma característica de glicose que
notifica uma mensagem de tempo real de 14~\textit{bytes}
(\textit{opcode}, \textit{status}, número de sequência de 32~\textit{bits},
\textit{timestamp} de 32~\textit{bits}, glicose de 12~\textit{bits}, um
\textit{byte} de estado e um \textit{byte} de tendência), anunciando
\texttt{FEBC} e o nome \texttt{DXCM01}. Não há J-PAKE, \textit{backfill} de
24\,h nem fluxo de calibração. O multiplicador de velocidade, a reprodução de
CSV e o modelo de falha PISA continuam todos valendo, pois moldam o valor de
glicose antes de a \texttt{comm\_thread} escolher o perfil. O ponto que esta
subseção sustenta é modesto e verificável: o firmware pode transmitir por
\emph{dois} formatos de fio distintos --- o CGMS da SIG e este segundo formato
não-SIG --- e a aplicação decodifica ambos, ou seja, o cliente não está preso
a um único protocolo. A validação desse ponto é o teste \textbf{S16-01} da
suíte E2E (força SIG, troca para o perfil Dexcom, volta para SIG, reconectando
a cada passo; os dois fluxos decodificam para valores de glicose plausíveis).
A mensagem de 14~\textit{bytes} está tabulada byte a byte em
\texttt{PROTOCOL\_SPEC.md}. Um transmissor Dexcom real de consumidor
\emph{não} lê este fluxo (esperaria o \textit{handshake} de autenticação e o
\textit{layout} EGV real); tornar a placa legível por um aplicativo de terceiro
que fale com Dexcom (via desafio AES do G6 + formato EGV real) é uma extensão
de escopo maior, rastreada à parte e não coberta por esta versão.

\subsection{Modelo de falha PISA}
\label{sec:impl-pisa}

O evento de atenuação do sensor induzida por pressão (\textit{Pressure-Induced
Sensor Attenuation}, PISA), cuja origem clínica --- a queda transitória de
perfusão no tecido em torno do eletrodo, que produz um falso valor baixo de
glicose --- foi descrita na fundamentação teórica, é reproduzido no simulador
como um evento pontual (característica \texttt{5b2c0013}, com carga útil
\texttt{\{u16 duration\_min; f32 depth\_frac\}}), na mesma classe
não-persistente e não-reiniciante dos eventos de refeição e exercício
instantâneos.

Enquanto um episódio de duração $T$ está ativo, a leitura final do sensor é
multiplicada pelo fator de atenuação da Equação~\ref{eq:pisa}:

\begin{equation}
a(t_{el}) = 1 - d\,\operatorname{sen}\!\left(\pi\,\frac{t_{el}}{T}\right)
\label{eq:pisa}
\end{equation}

Onde $T$ é a duração do episódio, em minutos; $t_{el} = T - t_{rem}$ é o
tempo decorrido desde o início do episódio, sendo $t_{rem}$ o tempo restante;
$d \in [0,1]$ é a profundidade de pico, recebida no campo \texttt{depth\_frac}
da característica; e $a(t_{el})$ é o fator multiplicativo aplicado à leitura.
O fator vale $1$ (nenhum efeito) no início e no fim do episódio e $1 - d$ no
seu ponto médio, descrevendo uma queda suave e simétrica; episódios
concorrentes se compõem multiplicativamente. Ele é aplicado \emph{depois} do
modelo de ruído do sensor e \emph{também} sobre a fonte de dados CSV, porque
o PISA é um artefato do transdutor e não uma mudança na glicose subjacente:
o valor verdadeiro de glicose no registro interno do firmware permanece
intacto, apenas a leitura reportada cai. O motor local da aplicação espelha o
mesmo fator, e o intervalo de tempo afetado é sombreado no gráfico.

Trata-se de uma aproximação fenomenológica, projetada para o simulador ---
dois parâmetros (profundidade e duração), com início, vale e recuperação
qualitativamente corretos ---, e não de um modelo publicado de atenuação por
pressão. O fenômeno clínico (\textit{compression low} / artefato por
compressão) e sua caracterização --- queda transitória e reversível da
leitura, tipicamente durante o sono ao deitar sobre o sensor, sem hipoglicemia
real --- estão descritos em Facchinetti et al., ``Modeling transient
disconnections and compression artifacts of continuous glucose sensors''
(\textit{Diabetes Technology \& Therapeutics}, 18(4):264--272, 2016); a forma
de meio-seno da Equação~\ref{eq:pisa} é uma escolha de projeto deste trabalho,
não a daquele artigo.

\subsection{Arquitetura multi-sensor}
\label{sec:multi-sensor}

\subsubsection{$N$ \textit{slots} independentes sobre uma placa}
\label{sec:multi-sensor-slots}

Uma opção de tempo de compilação, \texttt{CONFIG\_APP\_SENSOR\_COUNT} (de 1 a
4, padrão 4), define quantos sensores CGM totalmente independentes a mesma
placa simula. Cada \textit{slot} tem a sua própria estrutura de configuração
persistida: um modelo fisiológico com parâmetros, um modelo de ruído de
sensor e uma agenda diária recorrente de refeição e exercício --- \emph{ou}
uma trilha CSV enviada ---, e a combinação é livre (por exemplo, três
\textit{slots} de CSV e um de modelo). Cada \textit{slot} recebe ainda a sua
própria identidade BLE, o seu próprio conjunto de anúncio e a sua própria
instância do serviço GATT CGMS; a identidade~0 usa o endereço de fábrica e as
identidades $1$ a $N-1$ usam endereços aleatórios estáticos e estáveis, com
nome anunciado ``Nordic Glucose Sensor'' seguido do índice.

O que é \emph{compartilhado} entre todos os \textit{slots} é um único relógio
de simulação, um único multiplicador de velocidade e um único estado de
execução --- um só \emph{Start/Stop} comanda os quatro. Uma característica
\textit{Sensor select} (\texttt{5b2c0015}) funciona como um cursor por
conexão: seleciona qual \textit{slot} as leituras e escritas por sensor
seguintes (pessoa, sensor, fonte de dados, eventos, envio de CSV) endereçam,
o que permite ao serviço de configuração, de instância única, atender os $N$
\textit{slots} sem acrescentar um \textit{byte} de índice a cada mensagem. A
\texttt{comm\_thread} envia uma \textit{CGM Measurement} por \textit{slot} a
cada ciclo, além de uma notificação de \textit{Food/Exercise Status} por
\textit{slot}, esta última com um \textit{byte} de \textit{slot} no início.

No lado da aplicação, o usuário conecta ao endereço de cada identidade; cada
conexão é uma \texttt{BleSession} e cada sessão mostra apenas as medições do
seu próprio \textit{slot}, demultiplexadas pelo dígito final do nome
anunciado. Conectar às quatro identidades resulta em quatro linhas na árvore.
O caso $N = 1$ é a versão original de sensor único e é idêntica a ela no
enlace.

\subsubsection{Limitação de pareamento no Windows}
\label{sec:multi-sensor-windows}

Durante a validação em hardware constatou-se que o Windows conclui o
pareamento \textit{LE Secure Connections} contra a identidade BLE~0 (a de
fábrica) e realiza GATT criptografado com ela normalmente, mas \emph{aborta}
o mesmo \textit{handshake} contra as identidades não-padrão da placa. A placa
registra a mensagem \texttt{Security failed \ldots\ level 1 err 2}
(\texttt{BT\_SECURITY\_ERR\_AUTH\_REQUIREMENT}) e o Windows derruba o enlace
imediatamente (razão HCI \texttt{0x13}). As mitigações testadas em hardware e
que \emph{não} resolveram para as identidades $1$ a $N-1$ foram: desligar a
privacidade BLE, para que essas identidades usem endereços estáticos em vez
de endereços privados resolvíveis; reservar uma posição de vínculo por
identidade; e fazer a placa iniciar a segurança logo após a conexão.
\guilherme{Confirmar se essa limitação é do Windows testado (versão) ou geral; vale um teste rápido em Linux/Android para a monografia.}

A solução adotada é uma opção de compilação
(\texttt{CONFIG\_APP\_CGMS\_NO\_AUTH}, ativa por padrão quando
\texttt{CONFIG\_APP\_SENSOR\_COUNT} $> 1$ e inativa no \textit{build} de
sensor único): com ela, as características do serviço CGMS passam a exigir
apenas leitura e escrita simples, sem enlace autenticado, e todas as $N$
identidades transmitem sem qualquer pareamento. Isso é aceitável para um
simulador de bancada --- dados sintéticos, uso em laboratório, sem informação
pessoal ---, mas representa uma redução real de segurança: o enlace fica sem
criptografia, sem autenticação (qualquer central próxima pode conectar e,
pelo serviço de configuração, trocar modelos, subir um CSV ou parar a
simulação), sem proteção contra um intermediário que altere valores em
trânsito, e em desacordo com o perfil da SIG, que pressupõe enlace
autenticado para essas características. O \textit{build} de sensor único
mantém o pareamento e a criptografia normais. Caso um enlace multi-sensor
pareado e criptografado venha a ser necessário, os caminhos práticos são usar
uma pilha BLE de \textit{host} fora do Windows (por exemplo, um \textit{dongle}
nRF52840 atuando como central) ou placas físicas separadas.
